\begin{abstract}
在线策略蒸馏（On-Policy Distillation，OPD）通过学生自生成轨迹上的密集反馈，将教师的推理行为迁移给学生。常见的sampled-token实现为每个生成token分配单独的反馈信号，但一个局部推理决策可能由多个token共同表达。在短片段内共享反馈，能否改善这一学习过程？本文提出Block3 Mean，将连续三个token的师生对数概率差取平均，并将共享信号与联合block重要性比率、按有效长度加权的目标结合。该方法复用已有教师分数，不需要额外奖励模型或过程标注。我们分析其更新动力学，并在不同师生对和保存权重上评测数学推理能力。在历史Qwen3-1.7B对照的Step200，Block3相对sampled-token OPD的Avg@8在MATH500和AMC23上分别提高8.90和5.42个百分点；官方Qwen3-8B到1.7B指令版的Step200也在四项benchmark上获得正向Avg@8差值，而其他配对或保存点可能更有利于逐token更新。我们从局部信用共享、梯度尺度和裁剪几何解释这些比较所涉及的优化变化，并提供区分数值稳定性与生成质量的诊断。这些单seed实验与数学分析刻画了反馈粒度这一推理蒸馏设计选择，并指出值得进一步受控检验的条件。

\end{abstract}

\begin{figure}[htbp]
\centering
\includegraphics[width=\linewidth]{study_overview.png}
\caption{\textbf{从逐token反馈到block级更新。}学生生成响应，冻结教师在相同前缀上打分。Token OPD让每个token保留自己的反馈；Block3 Mean对三个相邻token的反馈求均值并共享，同时使用联合block ratio和按有效长度加权的block loss。上下两行是独立训练的对照方案，不是相加的两个目标。教师反馈与旧策略概率均停止梯度，仅更新当前学生。图示完整三token块；不足三token的尾块按有效长度$n_j$处理。两种方法均先采样完整轨迹再更新。}
\label{fig:overview}
\end{figure}

\section{引言}
从能力较强的教师迁移推理行为，是提升语言模型能力的重要途径。
在线策略蒸馏（On-Policy Distillation，OPD）让学生自主生成响应，并在其实际访问的前缀上接受教师反馈\citep{gkd,minillm}。
这种方式把监督带到学习者自身诱导的状态上，缓解只从专家示范学习时可能出现的分布错配\citep{dagger}。
教师概率还能够在整个响应中提供反馈，而不只在最终答案处给出信号。

其中仍有一个重要选择：应当如何将反馈分配给生成文本？
在sampled-token OPD中，每个token使用自身的师生对数概率差作为反馈\citep{revisiting}。
这种监督是密集的，但信号密度本身并不能决定信用分配应采用什么单位。
因此，我们提出问题：\textbf{由相邻token构成的短block，能否成为OPD中有效的反馈单位？}

我们的直觉是，一个局部推理决策的组成部分，例如短数学表达式或推理过渡，可能跨越多个token。
逐token概率差既反映背后的选择，也可能受到具体文本表达方式的影响。
共享局部信号有望降低对这些波动的敏感性，同时保留比整条响应奖励更细的粒度。
本文将这一点作为待检验的假设：固定token块是简单可操作的单位，并非经过标注的语义推理步骤。

基于这一视角，我们提出\method{}（图~\ref{fig:overview}）。
学生采样响应，冻结教师在相同前缀上打分；随后将响应划分为不重叠的三token块，
对块内蒸馏advantage求平均，并将共享advantage用于联合block重要性比率和按有效长度加权的目标。
这一方法复用已有教师分数，不需要额外奖励模型或过程标注。训练仍然在每次更新前收集完整轨迹。

我们在MATH500、AIME24、AIME25和AMC23上进行评测，逐benchmark报告结果。
每个匹配对照内，两种方法采用相同初始化、数据、提示和评测设置。
历史Qwen3-1.7B对照的Step200在MATH500和AMC23上的Avg@8分别提升8.90和5.42个百分点；
官方Qwen3-8B到1.7B指令版的Step200则在四项benchmark上均获得正向差值。
其他配对或保存点也存在更有利于Token的结果，因此师生条件与训练动力学是评测的重要组成部分。
这些结果按原有协议分别保留，不混合为独立重复实验。

本文有三项贡献。首先，我们构建结合局部advantage共享与联合重要性比率的block级OPD更新。
其次，我们推导其中的跨token信用项与梯度尺度，并给出平均操作改善局部信号估计的条件。
最后，我们提供逐checkpoint对照及覆盖熵、token重合、信用重分配和生成行为的轨迹诊断。
这些分析将方法的观测表现与其优化结构联系起来；单seed结果及因果归因的边界见第\ref{sec:results}和\ref{sec:limits}节。

\section{相关工作}
\label{sec:related}
GKD使用学生生成序列，并允许不同散度，因此on-policy采样并不唯一要求reverse KL\citep{gkd}。
MiniLLM研究语言生成中的reverse-KL蒸馏\citep{minillm}。
本文基线沿用Revisiting OPD的sampled-token实现；该研究分析局部蒸馏信号的不足与实用修复\citep{revisiting}。
我们不将这个基线等同于完整序列KL梯度的无偏估计。

Rethinking OPD研究师生兼容性及训练动力学，为本文监控熵、top-token重合和位置变化提供依据\citep{rethinking}。
这些是诊断工作的前例，不代表我们已复现该论文的全部最佳算法配置。
Blockwise policy-drift gating聚合学生新旧log-probability变化，构造停止梯度、均值归一化的权重，
但保留原有token目标\citep{bpdg}；本文则聚合教师反馈并改变PPO ratio本身。
FiRe-OPD过滤轨迹并对token重加权，TIP选择有信息量的token监督\citep{fireopd,tip}。
R2-OPD根据独立估计的推理进展过滤span\citep{r2opd}；GSPO采用长度归一化的序列ratio，不同于本文token ratio的乘积\citep{gspo}。
这些前例不支持泛称本文首次提出blockwise、选择性或多步OPD；我们也未直接评测它们的最佳方案。

\section{局部信用分配}
\label{sec:method}
\paragraph{采样局部反馈。}
给定prompt $x$，旧学生$\pi_{\mathrm{old}}$生成$y=(y_1,\ldots,y_T)$。
在前缀$h_t=(x,y_{<t})$上定义
\begin{equation}
 a_t=\log\pi_{\mathrm{T}}(y_t\mid h_t)-\log\pi_{\mathrm{old}}(y_t\mid h_t),
 \quad r_t(\theta)=\frac{\pi_\theta(y_t\mid h_t)}{\pi_{\mathrm{old}}(y_t\mid h_t)}.
 \label{eq:token}
\end{equation}
对$a_t$和旧概率停止梯度。这里的advantage是蒸馏log ratio，不是经过验证的任务正确性增量。
教师评分仍然自回归，且可在一次批量前向中评估已生成序列；将token分组不意味着无需计算各自的条件概率。

令$\mathcal C(r,a)$表示PPO裁剪代理目标\citep{ppo}，继承的负advantage dual-clip规则见附录\ref{app:proofs}。
忽略独立配置的正则项，基线在有效响应token上平均$-\mathcal C(r_t,\sg(a_t))$。
为便于说明，以下使用含$T$个有效token的单条响应；实现中mask padding，并在每个micro-batch的有效条目上归一化。

\paragraph{实际block干预。}
将连续响应tensor位置划为不重叠、宽度为$k$的block $B_j$，不跨mask空洞重新打包。
下式块内求和只选有效位置，$n_j$为其数量，$T=\sum_j n_j$，空块权重为零。
省略实现中分母的$10^{-8}$稳定项后，理想实数归一化为
\begin{align}
 \bar a_j&=\frac{1}{\max(1,n_j)}\sum_{t\in B_j}a_t,
 &R_j(\theta)&=\exp\left(\sum_{t\in B_j}\log r_t(\theta)\right),\label{eq:block}\\
 \mathcal L_{\mathrm{block}}&=-\frac1T\sum_{j:n_j>0}n_j\mathcal C(R_j,\sg(\bar a_j)),
 &\mathcal L_{\mathrm{token}}&=-\frac1T\sum_{t=1}^T\mathcal C(r_t,\sg(a_t)).\label{eq:loss}
\end{align}
\method{}取$k=3$。$R_j=\prod_{t\in B_j}r_t$不是token ratio的平均值。
代码的block mask权重为$n_j/k$，分母为这些分数权重之和；式\ref{eq:loss}是省略稳定项后等价的token计数形式。
完整等长块退化为按block平均；含2个有效token的尾块权重是$2/3$，不是1。
该目标作用于采样block，不枚举$|V|^k$联合分布。
两组均先生成完整轨迹batch再更新，不是每3个token更新一次优化器。

\begin{figure}[t]
\centering
\includegraphics[width=.94\linewidth]{credit_assignment_concept.png}
\caption{反馈共享示例，非实测rollout。$(2,-1,2)$平均后为$(1,1,1)$，中间token的局部信号翻转。
实际方法还改变ratio和损失归一化。局部符号翻转本身不证明任务层面的信用分配错误。}
\label{fig:credit}
\end{figure}

\paragraph{交叉信用项与尺度。}
考虑等长块、$T=kJ$，并在旧策略处计算，此时所有ratio为1且裁剪不生效。
记$g_t=\nabla_\theta\log\pi_\theta(y_t\mid h_t)$，策略梯度部分满足
\begin{align}
 \nabla\mathcal L_{\mathrm{token}}&=-\frac1T\sum_t a_tg_t,\nonumber\\
 \nabla\mathcal L_{\mathrm{block}}&=-\frac1T\sum_j\sum_{i\in B_j}\sum_{t\in B_j}a_i g_t
 =\nabla\mathcal L_{\mathrm{token}}-\frac1T\sum_j\sum_{\substack{i,t\in B_j\\i\ne t}}a_i g_t.
 \label{eq:cross}
\end{align}
额外项将token $i$的信号传给token $t$的score梯度。
它可能与基线同向，也可能相反；仅凭符号不能判断后续数学推理是否受益。

\begin{proposition}[相同反馈不意味着相同更新]
在$\theta=\theta_{\mathrm{old}}$、裁剪不生效、完整等长块大小为$k$且所有有效token满足$a_t=a$时，
两个策略损失均为$-a$，但$\nabla\mathcal L_{\mathrm{block}}=k\nabla\mathcal L_{\mathrm{token}}$。
\end{proposition}
将$\bar a_j=a$代入式\ref{eq:loss}即可得到该结论。
它描述的是策略损失导数，不意味着Adam参数更新、加入其他项后的总梯度或梯度裁剪后的更新也是$k$倍。
因此仅检查loss数值的单元测试可能漏掉实现差异。
离开旧策略后，联合ratio还改变裁剪区域：三个token ratio均为1.1时，block ratio为$1.1^3=1.331$。
token ratio仍可能在常见的0.2裁剪区间内，而block ratio已经越界，单独调整学习率一般无法消除此差异。

\paragraph{平均何时可能有帮助。}
为了将平均假设与上述具体实现分离，考虑一个假想干预：保持token ratio和token归一化，只将$a$换成$Wa$，
$W$为确定性矩阵，每行之和为1。仅作为局部信号模型，设$a=\mu+\varepsilon$、$\E\varepsilon=0$、
$\Cov(\varepsilon)=\Sigma$，则
\begin{equation}
 \E\|Wa-\mu\|_2^2=\|(W-I)\mu\|_2^2+\operatorname{tr}(W\Sigma W^\top).
 \label{eq:mse}
\end{equation}
只有偏差代价小于降噪收益时，平均才比原始信号的估计误差更低。
块内真实信号恒定、噪声独立同方差时，每位置方差变为$\sigma^2/k$。
真实信号不同或噪声高度相关时，这一优势可能消失。
式\ref{eq:mse}不是Block3或滑动窗口提升reward的定理：它研究信号估计，假定存在潜在分解，并固定梯度特征。
它也没有隔离历史方法的ratio及归一化变化。
附录\ref{app:proofs}给出推导，以及无条件平滑收益的反例。


\section{实验协议}
\label{sec:protocol}
\paragraph{证据单位。}
每个模型对、提示版本、数据池、grader和训练seed构成独立实验层。
同一层内，Token与Block3分别从同一个原始学生独立初始化，不从彼此权重继续训练。
每个checkpoint分别报告四个数学benchmark。
旧小规模评测和其他grader的结果只作为探索证据，不视为主对照的额外重复。
附录\ref{app:history}列出找到的全部实验系列，包括中止尝试和证据缺口。

\paragraph{匹配的4B Base训练。}
学生为Qwen3-4B-Base，而非同名指令模型\citep{qwen3}；教师为公开的Qwen3-4B-Base-GRPO，
具体身份见附录\ref{app:config}。两组复用既有DAPO衍生训练池和相同冻结训练实现。
本地parquet将17,917行题答序列重复100次，共1,791,700行，精确唯一题面为17,237个。
每步4个prompt，每个prompt采样8个回答，共200步；因此每组包含800次prompt出现、783个精确不同题面和6,400条轨迹。
训练seed为21，学习率$2\times10^{-6}$，每次更新1个PPO epoch，使用4张A100。
Step50、100、150、200均保存完整训练状态。
两组均实际遇到一次经字符串规范化后与MATH500重叠的题；附录\ref{app:config}报告训练池与实际暴露审计，以及单独的敏感性分析。

两组均使用temperature=1.0、top-$p=0.9$、2,048-token输入上限和16,384-token响应上限。
每个请求的生成seed由全局seed、step、数据身份、sample index和turn派生，不将同一seed重置后用于所有8个回答。
训练后审计核对全部6,400对输入身份与采样设置。
输入流相同不意味着后期输出相同：策略分化后生成不同on-policy数据是预期行为。
训练、生成和教师log-prob的micro-batch设置分别记录于附录\ref{app:config}。

\paragraph{提示控制。}
最新Base模型对照绕过ChatML，训练和评测使用完全相同的函数：
\begin{quote}\small
\texttt{[question]}\\
\texttt{Please solve the problem step by step and put the final answer in}\ \verb|\boxed{}.|\\
\texttt{Solution:}
\end{quote}
各部分间保留空行，\texttt{Solution:}后有最终换行。
显式关闭thinking，但允许普通逐步推理。合并checkpoint后检查真实输入token ID，而非仅比较文字。
EOS为151643，不增加其他stop-token ID。
旧Qwen和Llama实验保留并注明当时的提示协议，不把它们重新归类为最新统一completion设置。

\paragraph{评测与端点。}
新增官方指令版复用同一数据、优化器配置、seed21、200步及四个checkpoint。
提示保留相同boxed答案指令，但采用Qwen原生chat、显式关闭thinking、不加\texttt{Solution:}。
EOS为151645，停止集合为$\{151645,151643\}$，并增加原始学生对照。
附录\ref{app:instruct}说明模型revision、输入核验和恢复来源；它不是旧1.7B/公开GRPO分层的重复。

每个checkpoint评测MATH500（500题）、AIME24（30题）、AIME25（30题）和AMC23（83题），
共643题、5,144个回答。每题8次，采样seed为21至28，temperature=1.0、top-$p=0.9$，
两组均采用16,384-token上限。使用4个独立TP=1推理worker。
保留固定版本的历史数学grader，以及全部原始输出、输入token、采样seed和引擎停止原因。
评测显存比例为0.9，与训练rollout的0.6不同。CPU合并不修改源训练权重。

令$c_i\in\{0,\ldots,8\}$为第$i$题判对的回答数，则
\begin{equation}
 \avg=\frac{1}{N}\sum_{i=1}^N\frac{c_i}{8},\qquad
 \pass=\frac{1}{N}\sum_{i=1}^N\mathbf{1}[c_i>0].
 \label{eq:metrics}
\end{equation}
前者为平均采样准确率，后者为8次尝试至少成功一次的题目比例，均不是8个答案多数投票。
采用计划内最终Step200作为主要报告端点，这是回顾性的报告选择，而非预注册。中间checkpoint用于描述变化轨迹，不挑选某方法的最佳结果作为主结论。
不合并四个benchmark。报告的区间在各benchmark内部按成对题目重采样，条件固定为这些已训练checkpoint，
不代表不同训练seed之间的波动。

\paragraph{行为与数值诊断。}
保留学生/教师熵、带符号及绝对熵差、top-16重合率与重合概率质量、重合token advantage、
裁剪前梯度范数、策略梯度损失、响应长度和达到长度上限的比例。
位置热图汇总训练时实际访问的on-policy前缀，不是固定评测集。
block专用的符号分歧与advantage重分配指标作为补充。
缺失字面\verb|\boxed|标记和引擎length停止分别报告：二者都不等同于数学错误，缺少标记也只是特定格式失败指标。
精确定义和mask处理见附录\ref{app:diagnostics}。

\section{实验结果}
\label{sec:results}
\paragraph{不同师生对上的推理迁移。}
表~\ref{tab:historyendpoints}展示两组Step200 Avg@8差值为正的历史案例。
Qwen3-1.7B-Base学生与Rethinking OPD\citep{rethinking}发布的公开GRPO教师配对时，Block3的四项Avg@8均高于Token：
MATH500从45.85升至54.75，AIME24从5.42升至8.75，AIME25从2.50升至5.42，
AMC23从21.84升至27.26。这些观察推动我们进一步研究超出逐token更新的局部反馈共享。
DeepSeek/JustRL 1.5B对照的终点差值也为正，MATH500和AMC23的变化小于AIME。

这两组是正向案例展示，不是全部跨模型对照；其他历史师生对存在更有利于Token的结果。
附录\ref{app:history}表~\ref{tab:historyendpointsfull}保留完整四组终点对照及协议背景，
随后给出全部checkpoint结果。

\begin{table}[htbp]
\centering\small
\caption{选取的历史Step200 Avg@8正向案例（\%）。仅在每个学生分层内比较方法；包括负结果的完整对照见表~\ref{tab:historyendpointsfull}。}
\label{tab:historyendpoints}
\begin{tabular}{@{}llrrrr@{}}
\toprule
Student & Method & MATH500 & AIME24 & AIME25 & AMC23 \\
\midrule
Qwen 1.7B & Token & 45.85 & 5.42 & 2.50 & 21.84 \\
Qwen 1.7B & Block3 & 54.75 & 8.75 & 5.42 & 27.26 \\
\addlinespace
DeepSeek 1.5B & Token & 84.70 & 39.58 & 29.17 & 73.34 \\
DeepSeek 1.5B & Block3 & 84.75 & 43.75 & 30.83 & 73.80 \\
\addlinespace
\bottomrule\end{tabular}

\end{table}

\paragraph{官方Qwen3-8B到1.7B指令版。}
表~\ref{tab:qwen17endpoint}在原生非thinking chat、训推输入一致的设置下比较两种方法。
相对Token，Block3的Step200 Avg@8在MATH500、AIME24、AIME25和AMC23上分别高
$+0.10$、$+2.50$、$+0.83$和$+0.30$个百分点。
两种方法的四项Avg@8均超过原始学生。
方法差值的逐项95\%配对题目区间均包含零，因此正向点估计尚不能在该端点上明确区分统计优势。

Pass@8区分平均成功率与题目覆盖。
Block3相对Token的差值分别为$+0.20$、$-6.67$、$0.00$和$-1.20$个百分点，
且两种方法的MATH500覆盖率均低于原始学生的91.20\%。
因此，平均准确率提高可以与八次尝试覆盖率下降同时出现。
附录\ref{app:instruct}报告原始学生、全部八个训练权重和对应行为指标。
\begin{table}[t]
\centering\small
\caption{官方Qwen3-8B到Qwen3-1.7B指令版：Step200分数（\%）及Block3减Token差值（百分点）。}
\label{tab:qwen17endpoint}
\begin{tabular}{lrrrr}
\toprule
Benchmark & Original & Token & Block3 & $\Delta$ [95\% CI] \\
\midrule
\multicolumn{5}{c}{Avg@8} \\
MATH500 & 71.45 & 72.12 & 72.22 & +0.10 [-1.18, 1.38] \\
AIME24 & 10.00 & 14.17 & 16.67 & +2.50 [-2.50, 8.33] \\
AIME25 & 10.42 & 12.08 & 12.92 & +0.83 [-3.75, 5.83] \\
AMC23 & 40.51 & 43.07 & 43.37 & +0.30 [-3.16, 3.92] \\
\midrule
\multicolumn{5}{c}{Pass@8} \\
MATH500 & 91.20 & 89.60 & 89.80 & +0.20 [-2.00, 2.40] \\
AIME24 & 26.67 & 43.33 & 36.67 & -6.67 [-20.00, 6.67] \\
AIME25 & 23.33 & 30.00 & 30.00 & +0.00 [-16.67, 16.67] \\
AMC23 & 69.88 & 75.90 & 74.70 & -1.20 [-9.64, 7.23] \\
\bottomrule
\end{tabular}
\par\smallskip\footnotesize Original为本实验训练前的指令版学生，不是Base模型。方括号为10,000次配对题目重采样的逐点95\%区间；仅一个训练seed，未校正checkpoint或任务选择。
\end{table}


\paragraph{Qwen3-4B Base与公开GRPO教师。}
采用匹配completion提示的对照呈现不同的终点模式（表~\ref{tab:qwen4endpoint}）。
Step200时，Block3的Avg@8差值依次为$-0.52$、$-2.92$、$+0.42$和$-3.01$个百分点。
Token在三项任务上更高，但逐项95\%配对题目区间均跨零，既未确立优越性，也未证明等效。
AIME24中，Block3的Avg@8更低，却有多一道题至少答对一次，Pass@8为33.33\%，Token为30.00\%。
综合这些成对结果，信用粒度会影响学习轨迹，其收益与师生设置及所考察的指标有关。
% Generated from accepted results; requires booktabs.
\begin{table*}[t]
\centering
\scriptsize
\setlength{\tabcolsep}{4pt}
\caption{Qwen3-4B Step200 配对评测 (seed21; 队列 COMPLETE)。}
\label{tab:qwen4endpoint}
\begin{tabular}{lrrrrrr}
\toprule
 & \multicolumn{3}{c}{Avg@8} & \multicolumn{3}{c}{Pass@8} \\
\cmidrule(lr){2-4}\cmidrule(lr){5-7}
基准 & Block3 & Token & $\Delta$ [95\% CI] & Block3 & Token & $\Delta$ [95\% CI] \\
\midrule
MATH500 & 75.40 & 75.92 & -0.52 [-1.57, 0.55] & 90.60 & 91.00 & -0.40 [-2.40, 1.60] \\
AIME24 & 13.33 & 16.25 & -2.92 [-7.08, 0.83] & 33.33 & 30.00 & +3.33 [-6.67, 13.33] \\
AIME25 & 13.75 & 13.33 & +0.42 [-2.92, 3.75] & 30.00 & 30.00 & +0.00 [-13.33, 13.33] \\
AMC23 & 46.23 & 49.25 & -3.01 [-6.17, 0.30] & 74.70 & 77.11 & -2.41 [-10.84, 6.02] \\
\bottomrule
\end{tabular}
\par\smallskip\footnotesize 分数为百分比；$\Delta$为Block3减Token的百分点差。方括号为各基准内成对问题bootstrap的逐点95\%区间，非多训练种子CI；NA表示缺失。
\end{table*}


\IfFileExists{\qwenfigures/qwen4_accuracy.pdf}{%
\begin{figure}[t]
\centering
\includegraphics[width=\linewidth]{\qwenfigures/qwen4_accuracy.pdf}
\caption{当前匹配4B对照的四个保存权重。各benchmark独立报告Avg@8，使用独立纵轴范围，同一子图内两种方法共享范围；连线连接实测权重，不代表额外评测。每种方法仅一个训练seed。}
\label{fig:qwen4accuracy}
\end{figure}
}{}


\section{讨论与局限}
\label{sec:limits}
局部聚合存在偏差与方差的权衡，但历史实现同时改变了多个因素：共享奖励、概率比率相乘、归一化分母改变。
第\ref{sec:method}节的梯度计算说明，不能将所有涨点仅归因于相邻token语义一致。
关键后续对照应固定token ratio与token归一化，只改变平均算子，再独立改变ratio粒度及学习率尺度。
我们尚未完成这一完整析因实验，不将它描述为已有证据。

训练重复也是重要限制。单个训练seed即使评测数千个回答，仍然只有一条优化轨迹；
8次采样并不产生8次独立训练。AIME的题目分辨率尤其粗：多解决1题就会使Pass@8变化3.33个百分点。
窗口、checkpoint和先导权重经过适应性探索，不能将其中最高分当成选定方法优势的无偏估计。

提示诊断暴露了跨模型迁移实现的另一项限制：Base与指令模型对角色边界和thinking要求的反应未必相同。
首次优化更新之前已经存在的差异，不可能由block损失造成；但短提示probe也不能证明新协议在整个训练过程中消除了重复。
因此提示验收、优化器健康、生成质量和benchmark正确率分别检查。
位置熵有助于定位变化，但训练中on-policy状态分布也会变化，单凭热图不能识别因果机制。

最后，本文是与sampled-token基线的对照，不是对所有近期OPD算法的比较。
我们没有证明效率提升、教师计算减少或格式普遍正确。
4B completion表支持方法间差异，但没有匹配的原始学生结果；新增指令版具有这一对照，
却仍复用相同数据池，不构成独立数据集复现。更广泛的结论仍需要更多训练seed、组件消融和新任务分布。

\section{结论}
块级蒸馏改变的不只是教师汇总了多少个token。在本文研究的实现中，它同时改变信用分配位置、被裁剪的策略偏移以及策略损失导数的尺度。历史1.7B收益不足以支持普适结论：其他师生对出现混合结果或明显退化，而当前4B中间权重的收益随checkpoint反转，其终点在多数平均准确率benchmark上更有利于Token。实践上的要点是同时审计精确目标与生成轨迹。局部advantage共享仍然是可检验的假设；将其与ratio和归一化变化分离，是下一步应当完成的实验，而不是这些结果已经证明的收益。

\FloatBarrier
\label{mainend}

\section*{可复现性声明}
损失与推导见第\ref{sec:method}节和附录\ref{app:proofs}，两组成对协议见附录\ref{app:config}和\ref{app:instruct}，
历史版本与证据缺口见附录\ref{app:history}。
两组匹配对照均保留完整状态、全部轨迹和输入核验；表格与图由带来源的摘要生成。
这不意味着所有早期先导实验都有相同留存，也不意味着已公开匿名实验材料。

\section*{伦理声明}
本文使用数学推理benchmark与预训练语言模型，没有进行人类受试者数据收集。
更好地模仿教师本身不证明真实性或安全性。
我们保留负结果，不将格式成功、有限梯度或特定benchmark增益泛化为通用能力结论。
分发材料必须遵守上游模型与数据许可，并排除凭证、可识别身份的基础设施路径以及私有权重。

\section*{AI使用声明}
本研究使用生成式AI助手讨论和细化假设，辅助实验设计、数学论证及方法实现，调试和监控实验，
检索与总结相关文献，处理证据及解释结果，并起草、修订和翻译论文与图表。
学生和教师模型还生成构成方法本身的on-policy训练轨迹。
自动测试和材料审计检查特定实现与数据契约，但不能替代独立复现，也不能替代作者对最终投稿的责任。
本文不宣称研究、证明、写作或翻译未使用AI辅助。
